\chapter{What's New in Bijou 7.1.0}\label{record-282}

\noindent{\footnotesize\color{BookMuted}\path{docs/releases/7.1.0/whats-new.md}\par}\smallskip

Bijou 7.1.0 is a focused post-V7 minor release. It does not try to become the full V8 Runtime Graph product. Instead, it makes the first portable UI proof chain inspectable, adds a real DOGFOOD GraphQL block fixture, tightens theme and terminal-rendering behavior, and gives release prep better executable guardrails.

\section{Runtime Graph And Scene IR Proofs}\label{record-282-runtime-graph-and-scene-ir-proofs}

The core package now exposes the first \texttt{u\allowbreak{}i\allowbreak{}-\allowbreak{}s\allowbreak{}c\allowbreak{}e\allowbreak{}n\allowbreak{}e\allowbreak{}-\allowbreak{}i\allowbreak{}r\allowbreak{}/\allowbreak{}1} runtime contract for portable Bijou UI proof. The seed includes deterministic scene hashing, structural receipts, node/action/binding/token/i18n validation, visible-cell source-map facts, terminal lowering receipts, and lower modes for inspecting node, token, and localization state.

This is the first committed step toward a product-grade Runtime Graph lane: a Bijou surface can be described semantically, lowered deterministically, and verified without treating screenshots as the only evidence.

\section{GraphQL-Authored Blocks}\label{record-282-graphql-authored-blocks}

Bijou now has a constrained GraphQL authoring path for block artifacts:

% Presentation of original arrow notation. See flow-manifest.json.
\begin{center}
\begin{tikzpicture}[book flow node/.style={draw=BookHeading,line width=.65pt,fill=BookPanel,rounded corners=1.5pt,align=center,inner sep=5pt,minimum height=10mm,font=\fontsize{9}{11}\selectfont,execute at begin node={\hyphenpenalty=10000\exhyphenpenalty=10000}},book flow edge/.style={-{Latex[length=2mm]},draw=BookHeading,line width=.7pt}]
\node[book flow node,text width=58mm] (n0) at (0,0.0) {Graph\allowbreak{}QL SDL};
\node[book flow node,text width=58mm] (n1) at (0,-1.55) {bijou-block/\allowbreak{}1};
\node[book flow node,text width=58mm] (n2) at (0,-3.1) {ui-scene-ir/\allowbreak{}1};
\node[book flow node,text width=58mm] (n3) at (4,-3.1) {terminal Surface proof};
\node[book flow node,text width=58mm] (n4) at (4,-4.9) {graphql-bijou-block-debug/\allowbreak{}1 facts};
\draw[book flow edge] (n0) -- (n1);
\draw[book flow edge] (n1) -- (n2);
\draw[book flow edge] (n2) -- (n3);
\draw[book flow edge] (n2) -- (n4);
\end{tikzpicture}
\end{center}


\texttt{c\allowbreak{}o\allowbreak{}m\allowbreak{}p\allowbreak{}i\allowbreak{}l\allowbreak{}e\allowbreak{}G\allowbreak{}r\allowbreak{}a\allowbreak{}p\allowbreak{}h\allowbreak{}q\allowbreak{}l\allowbreak{}B\allowbreak{}i\allowbreak{}j\allowbreak{}o\allowbreak{}u\allowbreak{}B\allowbreak{}l\allowbreak{}o\allowbreak{}c\allowbreak{}k\allowbreak{}(\allowbreak{})} and \texttt{l\allowbreak{}o\allowbreak{}w\allowbreak{}e\allowbreak{}r\allowbreak{}B\allowbreak{}i\allowbreak{}j\allowbreak{}o\allowbreak{}u\allowbreak{}B\allowbreak{}l\allowbreak{}o\allowbreak{}c\allowbreak{}k\allowbreak{}T\allowbreak{}o\allowbreak{}U\allowbreak{}i\allowbreak{}S\allowbreak{}c\allowbreak{}e\allowbreak{}n\allowbreak{}e\allowbreak{}(\allowbreak{})} turn SDL into a stable \texttt{b\allowbreak{}i\allowbreak{}j\allowbreak{}o\allowbreak{}u\allowbreak{}-\allowbreak{}b\allowbreak{}l\allowbreak{}o\allowbreak{}c\allowbreak{}k\allowbreak{}/\allowbreak{}1} artifact and a portable scene. Group directives, grouped field membership, semantic source anchors, token refs, i18n fallback facts, actions, bindings, lower modes, and debug summaries are all part of the proof surface.

The release also checks in a real DOGFOOD NavigationListBlock GraphQL fixture, so the proof path is no longer only synthetic.

\section{Theme Tokens, Modes, And Inspector Work}\label{record-282-theme-tokens-modes-and-inspector-work}

The design-token lane is more concrete:

\begin{itemize}
\tightlist
\item
  \texttt{d\allowbreak{}e\allowbreak{}f\allowbreak{}i\allowbreak{}n\allowbreak{}e\allowbreak{}T\allowbreak{}h\allowbreak{}e\allowbreak{}m\allowbreak{}e\allowbreak{}(\allowbreak{})} creates immutable dark/light theme definitions.
\item
  \texttt{t\allowbreak{}o\allowbreak{}k\allowbreak{}e\allowbreak{}n\allowbreak{}R\allowbreak{}e\allowbreak{}f\allowbreak{}(\allowbreak{})} and \texttt{r\allowbreak{}e\allowbreak{}s\allowbreak{}o\allowbreak{}l\allowbreak{}v\allowbreak{}e\allowbreak{}T\allowbreak{}h\allowbreak{}e\allowbreak{}m\allowbreak{}e\allowbreak{}C\allowbreak{}o\allowbreak{}l\allowbreak{}o\allowbreak{}r\allowbreak{}R\allowbreak{}e\allowbreak{}f\allowbreak{}(\allowbreak{})} preserve semantic color facts.
\item
  \texttt{B\allowbreak{}I\allowbreak{}J\allowbreak{}O\allowbreak{}U\allowbreak{}\_\allowbreak{}D\allowbreak{}A\allowbreak{}R\allowbreak{}K} and \texttt{B\allowbreak{}I\allowbreak{}J\allowbreak{}O\allowbreak{}U\allowbreak{}\_\allowbreak{}L\allowbreak{}I\allowbreak{}G\allowbreak{}H\allowbreak{}T} are first-party presets.
\item
  \texttt{d\allowbreak{}e\allowbreak{}f\allowbreak{}i\allowbreak{}n\allowbreak{}e\allowbreak{}T\allowbreak{}h\allowbreak{}e\allowbreak{}m\allowbreak{}e\allowbreak{}S\allowbreak{}a\allowbreak{}f\allowbreak{}e\allowbreak{}P\allowbreak{}a\allowbreak{}i\allowbreak{}r\allowbreak{}s\allowbreak{}(\allowbreak{})} records reusable contrast pair matrices.
\item
  DOGFOOD uses a high-contrast shell theme family with explicit dark and light modes.
\item
  The Theme Lab page and \texttt{F\allowbreak{}1\allowbreak{}0} Theme Inspector expose token swatches, color reuse diagnostics, and safe-pair diagnostics.
\end{itemize}

These pieces make theme behavior easier to inspect and harder to fake with one-off color strings.

\section{Raster-To-Glyph And Title Surface Polish}\label{record-282-raster-to-glyph-and-title-surface-polish}

The TUI package now includes deterministic raw-RGBA-to-\texttt{S\allowbreak{}u\allowbreak{}r\allowbreak{}f\allowbreak{}a\allowbreak{}c\allowbreak{}e} rendering with ASCII, Braille, and quad modes. DOGFOOD uses that pipeline for release-title art, the V7 Launch Wake title surface, and image-to-glyph side-app previews.

The side app opened by \texttt{n\allowbreak{}p\allowbreak{}m\allowbreak{} \allowbreak{}r\allowbreak{}u\allowbreak{}n\allowbreak{} \allowbreak{}i\allowbreak{}m\allowbreak{}g} or \texttt{n\allowbreak{}p\allowbreak{}m\allowbreak{} \allowbreak{}r\allowbreak{}u\allowbreak{}n\allowbreak{} \allowbreak{}i\allowbreak{}m\allowbreak{}a\allowbreak{}g\allowbreak{}e\allowbreak{}-\allowbreak{}v\allowbreak{}i\allowbreak{}e\allowbreak{}w\allowbreak{}e\allowbreak{}r} can preview SVG, PNG, and PPM/PNM files, pan and zoom the view, switch glyph modes, preserve sampled colors, tune Braille contrast, and enable deterministic ordered dithering.

\section{DOGFOOD And Terminal Fixes}\label{record-282-dogfood-and-terminal-fixes}

Several release-review fixes landed in the same release window:

\begin{itemize}
\tightlist
\item
  terminal resize and display-resume repaint now invalidates the front buffer so blank cells repaint immediately after geometry changes
\item
  DOGFOOD Blocks navigation lets Up move from the first Block Preview child back to the parent row
\item
  ANSI foreground/background resets now honor scoped \texttt{3\allowbreak{}9} and \texttt{4\allowbreak{}9} resets without leaking colors into following labels
\item
  the quit confirmation modal no longer duplicates its controls
\item
  the DOGFOOD title perf HUD stays on the title screen and reports current frame timing
\end{itemize}

\section{Release Guardrails}\label{record-282-release-guardrails}

The release process is stricter and more replayable:

\begin{itemize}
\tightlist
\item
  \texttt{n\allowbreak{}p\allowbreak{}m\allowbreak{} \allowbreak{}r\allowbreak{}u\allowbreak{}n\allowbreak{} \allowbreak{}r\allowbreak{}e\allowbreak{}l\allowbreak{}e\allowbreak{}a\allowbreak{}s\allowbreak{}e\allowbreak{}:\allowbreak{}r\allowbreak{}e\allowbreak{}a\allowbreak{}d\allowbreak{}i\allowbreak{}n\allowbreak{}e\allowbreak{}s\allowbreak{}s\allowbreak{} \allowbreak{}-\allowbreak{}-\allowbreak{} \allowbreak{}-\allowbreak{}-\allowbreak{}m\allowbreak{}i\allowbreak{}l\allowbreak{}e\allowbreak{}s\allowbreak{}t\allowbreak{}o\allowbreak{}n\allowbreak{}e\allowbreak{} \allowbreak{}v\allowbreak{}7\allowbreak{}.\allowbreak{}1\allowbreak{}.\allowbreak{}0} prints a milestone report and blocks incomplete release prep before running the full local gauntlet
\item
  DOGFOOD i18n debt discovery now scans \texttt{e\allowbreak{}x\allowbreak{}a\allowbreak{}m\allowbreak{}p\allowbreak{}l\allowbreak{}e\allowbreak{}s\allowbreak{}/\allowbreak{}d\allowbreak{}o\allowbreak{}c\allowbreak{}s\allowbreak{}/\allowbreak{}*\allowbreak{}*\allowbreak{}/\allowbreak{}*\allowbreak{}.\allowbreak{}t\allowbreak{}s} modules by default, with explicit tooling-only exclusions
\item
  release evidence packets are now required for versioned releases
\item
  the roadmap, bearing, and release runbook now describe the forward-looking V7.1/V8/V9 release posture
\end{itemize}

The practical result is a release that is easier to inspect: users get new portable proof primitives and DOGFOOD fixture evidence, and maintainers get better release gates before tags are cut.

